\begin{center}
RULES OF PROCEDURE\end{center}


\begin{center}
CHAPTER I.\end{center}
\begin{center}
\emph{Definitions.}\end{center}


1. These rules shall be called the Constituent Assembly Rules. They shall come into force at once.



2. In these rules, unless the context otherwise requires—



(a) "Assembly" means the Constituent Assembly of India;

(b) "Chairman" means the person who for the time being presides over the Assembly or any of its Committees;

(c) "Member" means a member of the Assembly;

(cc)  ‘Minister’ means a Member of the Council of Ministers of the Governor-General of India;

(d) "President" means the person elected by the Assembly at its preliminary session as the Chairman under the provisions of the Statement and his successors in office.

(e) "Secretary" means the Secretary of the Assembly.

(f) "Statement" means the Statement of the Cabinet Mission to India and H.E. the Viceroy dated May 16, 1946.

\begin{center}
\end{center}
\begin{center}
CHAPTER II.\end{center}
\begin{center}
\emph{Ad}\emph{m}\emph{ission of Members and Vacation of Seats}\end{center}
\begin{center}
\end{center}
3. No member shall take his seat or vote in the Assembly until he has signed his name in the Register kept for the purpose at a meeting of the Assembly or, if the Assembly is not in session, in the presence of the President.



Provided that every Minister who is not a Member of the Assembly shall have the right to speak in and otherwise to take part in the proceedings of the Assembly and any Committee thereof of which he may be named a member, but shall not be entitled to vote.



4. (1) A member may resign his office by writing under his hand addressed to the President.



(2) On the acceptance of the resignation by the President the seat shall become vacant.



5. (1) When a vacancy occurs by reason of death, resignation or otherwise in the office of a member of the Assembly, the President shall notify the vacancy and call upon the constituency concerned to elect a person for the purpose of filling the vacancy.



(2) Upon the occurrence of the vacancy, the President shall ordinarily make a request in writing to the Speaker of the Provincial Legislative Assembly concerned, or, as the case may be, the President of the Coorg Legislative Council, for the election of a person, for the purpose of filling the vacancy as soon as may reasonably be practicable.



(3) Where the vacancy is in the office of a member elected by a Provincial Legislative Assembly, the seat shall be filled by election according to the principle of proportional representation with the single transferable vote by the members of the same community whether General, Muslim or Sikh, in the Provincial Legislative Assembly, as had elected him.



(4) A candidate for a vacancy shall be such as could be elected by the constituency concerned under the Statement.



(5) Only an Indian, that is to say, a person domiciled in any part of India which is participating or entitled to participate in this Assembly, may be nominated for election as a member of the Constituent Assembly.



(6) As soon as may be after the receipt of the request mentioned in sub-rule (2) the Speaker of the Provincial Legislative Assembly concerned—



(a) shall appoint by suitable notification a person to be the Returning Officer for the election and may also in like manner appoint any person who may, subject to the control of the Returning Officer, perform all or any of the functions of the Returning officer at any such election, and



(b) shall also appoint by suitable notification—



(i) a date, not later than fifteen days after the date of notification for the nomination of candidates;

(ii) a further date, not later than the third day after the first-mentioned date, for the scrutiny of nominations;

(iii) a further date, not later than two days after scrutiny, for withdrawal of his candidature by a candidate; and

(iv) a further date, not later than twenty one days from the date fixed for withdrawal on which a poll shall if necessary, be taken.



(6)A. The Speaker of the Provincial Legislative Assembly concerned shall, if a poll is taken, by suitable notification fix the hour at which the poll shall commence and the hour at which it shall close on the date fixed under sub-clause (iv) of clause (b) of sub-rule (6) and the place at which the poll shall be taken.



(7) The votes shall be given by ballot and in person:

Provided that when the Provincial Legislative Assembly is not in session, votes may, at the option of the voter, be given in person or by registered post, provided further that no votes shall be given by proxy.



(8) After the votes are duly counted, the result of the election shall be declared and reported to the President.



(9) Save as otherwise provided in these Rules, the elections to the Constituent Assembly, shall be held, \emph{mutatis mutandis}, in accordance with the rules and regulations for the time being in force in regard to elections held by the Provincial Legislative Assembly, and, where no such rules and regulations exist, in accordance with the standing orders that may be made in this behalf by the President of the Constituent Assembly, where any such rules or regulations exist, it shall be competent for the Speaker of the Provincial Legislative Assembly concerned to make, with the previous approval of the President, such modifications therein as may be necessary for the purposes of this sub-rule.



(10) The foregoing rules shall apply in relation to Coorg, subject to the following modifications, namely,—



(a) that for the "Provincial Legislative Assembly" there shall be substituted "the Coorg Legislative Council" and for the "Speaker of the Provincial Legislative Assembly" there shall be substituted the "President of the Coorg Legislative Council", and

(b) that instead of a section of the Provincial Legislature taking part in the election, the non-official members of the Coorg Legislative Council shall take part in it.



(11) The foregoing rules shall apply in relation to Delhi and Ajmer-Merwara subject to the following modifications, namely:



(a) that for the ‘the Provincial Legislative Assembly’ there shall be substituted ‘the Delhi Advisory Council or the Ajmer-Merwara Advisory Council, as the case may be; and for the ‘the Speaker of the Provincial Legislative Assembly’ there shall be substituted ‘the Chairman of the Delhi or Ajmer- Merwara Advisory Council as the case may be’.



(b) that instead of a section of the Provincial Legislature taking part in the election, the non-official members of the Delhi or the Ajmer-Merwara Advisory Council shall take part in it.



(12) If any vacancy occurs by reason of death, resignation, or otherwise in the office of a member representing Delhi or Ajmer-Merwara in the Constituent Assembly, the President shall notify the vacancy and shall call upon the Chief Commissioner of Delhi or Ajmer-Merwara as the case may be, to take steps to hold, a bye-election to fill the vacancy. The bye-elections shall be held, as nearly as may be, in accordance with the procedure prescribed by the Legislative Assembly Electoral Rules, as in force on August 1, 1947, for the election of a member to represent Delhi or, as the case may be, the Ajmer-Merwara constituency of the Indian Legislative Assembly



5-A. When a vacancy occurs by reason of death, resignation or otherwise in the office of a member of the Assembly representing an Indian State, the President shall notify the vacancy and make a request in writing to the Ruler of the Indian State concerned to proceed to fill the vacancy, as soon as may reasonably be practicable, by election or nomination, as the case may be. 



Provided that where the seat was filled previously by nomination, the Ruler may fill the vacancy by election.



5-B. In the case of a vacancy in the office of a member of the Assembly representing more than one Indian State, the President shall notify the vacancy and make a request in writing to the Rulers of the Indian States concerned to, proceed to fill the vacancy, as soon as may reasonably be practicable, by the same method as was applicable to the case of the outgoing member when he was chosen as a member of the Assembly.



\begin{center}
CHAPTER III.\end{center}
\begin{center}
\emph{The President.}\end{center}


6. (1) The President of the Assembly shall be elected by the Assembly from among its members.



(2) The President shall cease to hold office as such if he ceases to be a member of the Assembly.



(3) The President may resign office by writing under his hand addressed to the Secretary for communication to the Assembly. If the Assembly is in session, the resignation shall be read out to the members; if the Assembly is not in session, it shall be published in the \emph{Gazette of India}; upon being so read out or published, as the case may be, it shall become effective.



7. (1) When, owing to a vacancy in the office of President of the Assembly, the election of President becomes necessary, one of the Vice-Presidents appointed in this behalf by the Steering Committee, shall fix a date for the holding of the election and the Secretary shall send to every member notice of the date so fixed.



(2) At any time before noon on the day preceding the date so fixed, any member may nominate another member for election by delivering to the Secretary a nomination paper signed by himself as proposer and by a third member as seconder and stating—



(a) the name of the member nominated, and

(b) that the proposer has ascertained that such member is willing to serve as president if elected.



(3) Any member who has been nominated may withdraw his candidature at any time before the Assembly proceeds to hold the election.



(4) On the date fixed for an election, the Vice-President shall read out to the Assembly the names of the members who have been duly nominated and have not withdrawn their candidature together with those of their proposers and seconders, and, if there is only one much member, the Assembly shall proceed to elect the President by ballot.



(5) Where there are only two candidates for election, the candidate who obtains at the ballot the larger number of votes shall be declared elected. If they obtain an equal number of votes, the election shall be by the drawing of lots.



 (6) Where more than two candidates have been nominated and at the first ballot no candidate obtains more votes that the aggregate votes obtained by the other candidates, the candidate who has obtained the smallest number of votes shall be excluded from the election, and balloting shall proceed, the candidate obtaining the smallest number of votes at each ballot being excluded from the election, until one candidate obtains more votes that the remaining candidate or than the aggregate votes of the remaining candidates, as the case may be, and such candidate shall be declared elected.



(7) Where at any ballot any of three or more candidates obtain an equal number of votes and one of them has to be excluded from the election under sub-rule (6) the determination, as between the candidates whose votes are equal, of the candidate who is to be excluded shall be by the drawing of lots.



8. (1) The deliberations of the Assembly shall be presided over by the President, when he is present.



(2) The President shall be the guardian of the privileges of the Assembly, its spokesman and representative and its highest executive authority.



9. The President may, subject to such conditions as he may prescribe, delegate to a Vice-President such of his powers and duties as he may think fit.



\begin{center}
CHAPTER IV.\end{center}
\begin{center}
\emph{Vice-Presidents.}\end{center}


10. The Assembly shall have two Vice-Presidents who shall be elected by the Assembly from amongst its members in such manner as the President may prescribe.



11. The President shall fix the date and time for the nomination and election of each of the two Vice-Presidents to be elected under rule 11.



12. (1) A Vice-President shall cease to hold office as such if he ceases to be a member of the Assembly. A Vice-President may resign office by writing under his hand addressed to the President, and, on acceptance thereof by the President, the resignation shall become effective.



(2) Any vacancy in the office of a Vice-President of the Assembly shall be filled by an election by the Assembly from among its members at such time and in such manner as the President may prescribe.



13. In the absence of the President, such Vice-President as the President may determine shall preside over the Assembly.



14. If the President is absent and there is no Vice-President present to preside over the Assembly, the Assembly may choose any member to perform the duties of the Chairman.



\begin{center}
CHAPTER V.\end{center}
\begin{center}
\emph{Office of the Constituent Assembly.}\end{center}
\begin{center}
\end{center}
15. (1) The President shall be the head of the office of the Assembly.



(2) The Office shall consist of two branches, an Advisory Branch and an Administrative Branch.



(3) There shall be a Constitutional Adviser, to be appointed by the President, who shall be the head the Advisory Branch.



(4) There shall be a full-time Secretary, to be appointed by the President, who shall be the head of the Administrative Branch.



(5) There shall be under the Constitutional Adviser and the Secretary respectively such other officers and establishments as the President may, on the advice of the Staff and Finance Committee, determine.



(7) The President may, in consultation with the Government of any Governor’s Province, appoint a Secretary for the Province, who shall rank as a Deputy Secretary in the Office of the Assembly.



(8) Any member of the Assembly who becomes the Secretary of a Deputy Secretary under this rule shall thereupon be deemed to have vacated his office as member.



(9) The President shall exercise in respect of the Office of the Assembly all powers of appointment, control and discipline, provided that he may delegate to any officer such of these powers as he thinks fit and subject to such conditions as he may prescribe.



16. It shall be among the duties of the Secretary—



(a) to administer the funds placed at the disposal of the Assembly in accordance with the provisions in the budget accepted by the Assembly;

(b) to keep and maintain the records relating to the business of the Assembly; and

(c) to make arrangements for meetings of the Assembly and its Committees.



\begin{center}
CHAPTER VI.\end{center}
\begin{center}
\emph{Business of the Assembly.}\end{center}


17 (1) The business of the Assembly means the business conducted—



(i) in the Assembly itself;

(ii) in the Committee of the whole Assembly;

(iv) in the Advisory Committee referred to in paragraph 20 of the Statement;

(v) in the Steering Committee;

(vi) in the Staff and Finance Committee;

(vii) in the Credentials Committee;

(viii) in the House Committee; and

(ix) in such other committees or sub-committees as may be set up by the Assembly or the Advisory Committee or other committees.



(2) The business of the Assembly shall be conducted in New Delhi unless the Assembly resolves otherwise.



(3) The conduct of the business of the Assembly and the procedure thereof shall be regulated by the rules and standing orders and the resolutions of the Assembly and by the rulings given from time to time by the Chairman.



18. The Assembly shall not be dissolved except by a resolution assented to by at least two-thirds of the total number of members of the Assembly.



19. The Assembly shall sit on such dates as the President, having regard to the state of business of the Assembly, may from time to time direct:



Provided that the President shall not adjourn the session for more than three days at a time except with the consent of the Assembly:



Provided further that the Chairman may adjourn the session to the next working day.



20. The Assembly may resolve itself into a Committee of the whole Assembly.



21. Meetings of the Assembly shall commence at 11 A.M., except as otherwise resolved by the Assembly or directed by the Chairman.



22. (1) The presence of at least one-third of the whole number of members shall be necessary to constitute a meeting of the Assembly or any of its committees.



(2) If the Chairman, on a count being demanded by a member at any time during a meeting, ascertains that one-third of the whole number of members are not present, he shall adjourn the Assembly or the committee, as the case may be, for fifteen minutes, and if on a fresh count being taken after that period it is found that there is still no quorum, he shall adjourn the Assembly or the committee as the case may be, till the next day on which it ordinary sits.



23. (1) A list of business for the day shall be prepared by the Secretary and a copy thereof shall be supplied for the use of every member before the commencement of the business of the day. The business thus prepared shall be called the “Orders of the Day”. 



(2) Save as otherwise provided in these rules, no business, not included in the Orders of the Day, shall be transacted at any meeting without the leave of the Chairman.



24. The business of the Assembly shall be brought before it or its committees by means of—



(a) a motion;

(b) a report of a committee; 

(c) an amendment to a motion or an amendment to an amendment.



25. (1) Unless otherwise directed by the Chairman, notice of every motion, accompanied by a copy of the motion shall be given at least three clear days before the day on which the motion is to be moved in the Assembly.



(2) Every notice required by these rules shall be given in writing addressed to the Secretary and signed by the member giving notice and shall be left at the Notice Office, which shall be open for the purpose between the hours of 11 A.M. and 5 P.M. except on Sundays and other public holidays.



(3) Notices left when the office is closed shall be treated as given on the next open day.



(4) Where three clear days’ notice of a motion has been given, the Secretary shall send a copy of the motion to the members at least two clear days before the day on which it is to be moved; and in other cases, he shall send a copy to them as soon as possible after notice has been received.



(5) No notice shall be required—



(a) for a motion for an adjournment of the Assembly;

(b) for a motion for adjournment of the consideration of the motion which is under discussion;

(c) for a motion for reference back to a committee;

(d) for a motion that the Assembly do resolve itself into committee of the whole; or

(e) for a motion which, in the opinion of the Chairman, is of urgent and vital importance.



26. There shall be no motion for an adjournment of the Assembly for the purpose of discussing any matter not included in the Orders of the Day, or not connected with the work of the Assembly.



27. The members shall sit in such order as the President may direct.



28. A member desiring to make any observation on any matter before the Assembly shall only speak when called upon to do so, by, the Chairman and shall rise when he speaks, except as permitted by the Chairman. If at any time, the Chairman rises, the member shall take his seat and cease speaking.



29. (1) In this Assembly, business shall be transacted in Hindustani (Hindi and Urdu) or English, provided that the Chairman may permit any member who cannot adequately express himself in either language to address the Assembly in his mother tongue. The Chairman will make arrangements for giving the Assembly, whenever he thinks fit, a summary of the speech in a language other than that used by the member and such summary shall be included in the record of the proceedings of the Assembly.



(2) The official records of the proceedings of the Assembly shall be kept in Hindustani (both Hindi and Urdu) and English.



30. (1) A matter requiring the decision of the Assembly shall be brought forward by means of a question put by the Chairman.



(2) In all matters requiring to be decided by the members of the Assembly, the Chairman shall exercise a vote only in the case of an equality of votes.



(4) Votes may be taken by voices or division and shall be taken by division if any member so desires.



(5) The Chairman shall determine the method of taking vote by division.



(6) The result of a division shall be announced by the Chairman and shall not be challenged.



31. (1) An amendment must be relevant to the motion to which it is proposed.



(2) An amendment may not be moved which has the effect of being the negative of the original motion.



(3) Except as permitted by the Chairman—



(a) notice of any amendment to a motion must be given at least one clear day before the motion is to be moved in the Assembly, and

(b) notice of any amendment to an amendment must be given before the Assembly meets for the day on which the motion is to be moved.



(4) The Chairman may disallow any amendment which he considers to be frivolous or dilatory.



(5) The Chairman may put amendments to the vote in any order he may choose.



32. No question which has once been decided by the Assembly shall be reopened except with the consent of at least one-fourth of the members present and voting.



33. Any time after a motion has been made, any member may move “that the question be now put” and unless it appears to the Chairman that the motion is an infringement of the right of reasonable debate, the Chairman shall put the motion, “that the question be now put”; and if the motion is accepted, no further discussion on the original motion shall be permitted except for a reply by the member who made the original motion.



34. In all matters relating to procedure or the conduct of business of the Assembly, the decision of the Chairman shall be final.



35. It shall be the function of the Advisory Committee referred in paragraphs 19 and 20 of the Statement to initiate and consider proposals and to make a report to the Assembly upon fundamental rights, clauses for the protection of minorities and the administration of tribal and excluded areas; and the decisions of the Assembly upon such report shall be incorporated in the appropriate part of the Constitution.



36. The Chairman of the Assembly after having called the attention of the Assembly to the conduct of a member, who persists in irrelevance or in tedious repetition, either of his own arguments or of the arguments used by other members in debate, may direct him to discontinue his speech, and the member shall, thereupon, resume his seat.



37. (1) The Chairman shall preserve order and shall have all powers necessary for the purpose of enforcing his decisions on all points of order.



(2) The Chairman may, in the case of grave disorder arising in the Assembly, suspend any sitting for a period not exceeding three days.



38. (1) The admission of persons other than members to the Assembly Chamber and its galleries during the sittings of the Assembly shall be regulated in accordance with the orders of the Chairman:



Provided that when a meeting is held \emph{in camera} no persons other than members and the officers and staff on duty shall be admitted inside the Assembly Chamber or its galleries.



(2) The meetings of the Assembly may, in the discretion of the Chairman, be held \emph{in camera}.



(3) The proceedings of all committees shall be conducted \emph{in camera}.



(4) The Secretary shall cause full reports of the proceedings of the Assembly to be printed and supplied to all members:



Provided that, where any meeting is held \emph{in camera}, such report shall be marked confidential and for the personal use of the members only.



(5) Where any meeting of the Assembly is held \emph{in camera}, the Chairman may authorise a summary of the proceedings to be issued to the press.



\begin{center}
CHAPTER VII\end{center}
\begin{center}
\emph{Committees.}\end{center}


39. (1) A Steering Committee shall be set up for the duration of the Assembly and shall consist initially of eleven members (other than the President) to be elected by the Assembly in accordance with the principle of proportional representation by means of the single transferable vote.



(2) The Assembly may from time to time elect, in such manner as it may deem appropriate, eight additional members, of whom four shall be reserved for election from among the representatives of the Indian States.



(3) The President shall be an \emph{ex-officio} member of the Steering Committee and shall be its \emph{ex-officio} Chairman. The Committee may elect a Vice-Chairman from among its members to preside over the Committee in the absence of the President.



(4) The Secretary of the Assembly shall be \emph{ex-officio} Secretary of the Steering Committee.



(5) Casual vacancies in the Committee shall be filled as soon as possible after they occur by election by the Assembly in such manner as it may determine.



40. (1) The Committee shall—



(a) arrange the order of business for the day;

(b) group similar motions and amendments and secure, if possible, assent of the parties concerned to composite motions and amendments;

(c) act as a general liaison body between communities \emph{inter se} and between the President and any part of the Assembly; and

(d) deal with any other matter under the rules or referred to it by the Assembly or the President.



(2) The President may make standing orders for the conduct of the business of the Steering Committee.



41. (1) A Staff and Finance Committee shall be set up for the duration of the Assembly and shall consist of—



(a) the President, who shall be \emph{ex-officio }Chairman of the Committee,

(b) the two Vice-Presidents, and

(c) nine other members to be elected by the Assembly in accordance with the principle of proportional presentation by means of the single transferable vote.



(2) The Functions of the Committee shall be—



(a) to advise the President regarding the posts to be created in the Office of the Assembly, and the salaries and emoluments to be attached thereto;

(b) to recommend to the Assembly the allowances to be paid to the officers and members of the Assembly and its committees; and

(c) to frame a budget or supplementary budget for submission to the Assembly.



(3) The President may make standing orders for the conduct of the business of the Committee.



(4) The Staff and Finance Committee shall invite the Auditor-General to audit the accounts of the Assembly.



42. A Credentials Committee shall be set up for the duration of the Assembly for the purpose of dealing with all questions relating to the validity of the title of elected or other members.



43. (1) The Committee shall consist of five members who shall be elected by the Assembly.



(2) The Committee shall have power to co-opt additional members not exceeding two in number.



(3) Casual vacancies in the Committee shall be filled as soon as they occur by election by the Assembly or by co-option, as the case may be.



(4) The Committee shall elect its own Chairman.



(5) The President may make standing orders for the conduct of the business of the Committee.



44. (1) A House Committee shall be set up for the duration of the Assembly for the purpose of dealing with all questions relating to housing and accommodation of members, to provide for and to supervise over facilities for food, medical relief, entertainments, and reading room and library during their stay.



(2) The Committee shall consist of eleven members, who shall be elected by the Assembly, in the manner to be prescribed by the President.



(3) The Committee shall have power to co-opt additional members and to appoint sub-committees to deal with various items of their work.



45. (1) Any other committee may be set up by a motion in the Assembly.



(2) The members of every such committee shall, unless the motion by which the committee is set up otherwise provides, be elected according to the principle of proportional representation by means of the single transferable vote.



(3) Unless the motion by which a committee is set up otherwise provides, the committee shall appoint one of its members as Chairman who shall regulate the work of the committee.



46. Subject to the provision of these rules, the Secretary of the Assembly, shall be \emph{ex-officio} Secretary of every committee, unless the motion by which the committee is set up otherwise provides.



47. The motion by which a committee is to be set up may state the quorum necessary to constitute a meeting of the committee and may fix the time within which the committee shall present its report, if any.



48. The report of a committee shall be presented to the Assembly by the Chairman of the committee unless the committee otherwise provides, and a copy thereof shall be forwarded to the Office of the Assembly.



\begin{center}
CHAPTER VIII.\end{center}
\begin{center}
\emph{Budget.}\end{center}


49. (1) A statement of the estimated expenditure of the Assembly shall be prepared by the Staff and Finance Committee and placed before the Assembly for sanction.



(2) Supplementary statements may be similarly placed before the Assembly in accordance with the directions of the President.



\begin{center}
CHAPTER IX.\end{center}
\begin{center}
\emph{Salaries and Allowances.}\end{center}


50. (1) Allowances of the members shall be fixed by the Assembly on a motion approved by the Staff and Finance Committee. The Committee may make standing orders providing for special allowances to members in particular cases as well as for allowances to non-members engaged in the work of the Assembly.



(2) The salaries and allowances of the servants of the Government of India or any Provincial Government or any other authority whose services are placed at the disposal of the Assembly shall be such as may be agreed upon between the Government or other authority concerned and the President acting on the advice of the Staff and Finance Committee. The Salaries and allowances of all persons recruited directly shall be such as may be fixed by the President on the advice of the Staff and Finance Committee.



\begin{center}
CHAPTER X.\end{center}
\begin{center}
\emph{Doubts and Disputes as to Elections.}\end{center}


51. In this Chapter, unless there is anything repugnant in the subject or context,—



(a) “Candidate” means a person who has been nominated as a candidate at any election or who claims that he has been so nominated or that his nomination has been improperly refused’



(aa) ‘representative’ of any Indian State or States means the person who is chosen as a representative of such State or States in the Assembly in accordance with the provisions contained in the Schedule to these Rules.



(b) “Returned Candidate” means a candidate whose name has been published in the appropriate Official Gazette as a duly elected member of the Assembly and includes a candidate whose name has been reported by or on behalf of the Ruler or Rulers of any Indian State or States to the President in the manner provided in the Schedule to these rules as a duly chosen representative of such State or States.



52. No election shall be called in question except by an election petition presented in accordance with the provisions of this Chapter.



53. (1) An election petition against any returned candidate may be presented to the President by any candidate or elector on the ground of any irregularity or corrupt practice:.



(i) in the case of election to the Assembly held before the publication of these Rules, within fifteen days from the date on which these rules are published in the Gazette of India;



(ii) in the case of subsequent elections within thirty days from the date on which the results of the elections are published in the Gazette of India or in the Official Gazette of the Province concerned.



(2) An election petition shall be deemed to have been presented to the President when it is delivered to the President or to any officers appointed by him in this behalf:



(a) by a person making the petition; or

(b) by a person authorised in writing in this behalf by the person making the petition; or

(c) by registered post.



54. At the time of the presentation of the petition, or, in the case of any petition presented before the date on which these rules come into force, within fifteen days from the date on which these rules are published in the Gazette of India, the practitioner shall deposit the sum of rupees one thousand in cash or in Government promissory notes of equal value at the market rate of the day as security for the costs of the case.



55. (1) If the provisions of rule 54 or rule 55 are not complied with, the President shall dismiss the petition:



Provided that if the person making the petition satisfies the President that sufficient cause existed for his not presenting the petition within the period described in rule 54, the President shall have discretion to condone the failure to comply with that rule.

(2) If the petition is not dismissed under the preceding sub-rule, the President shall, if he is satisfied that sufficient grounds exist for such action, refer the petition to the Credentials Committee.



56. The Credentials Committee shall, with due despatch, inquire into the allegations made in the petition and, subject to the provisions of the next succeeding rule, submit a report to the President.



57. The Credentials Committee may, if they think fit, recommend to the President that an election tribunal be appointed to inquire into the petition.



58. Where such a recommendation has been made, the President shall appoint an election tribunal consisting of one or more than one person to inquire into the petition. 



59. On the conclusion of the inquiry, the election tribunal shall make a report to the President.



59-A. (1) The Credentials Committee or the Election Tribunal shall, for the purposes of an inquiry into an election petition, have power to summon and enforce the attendance of witnesses and to compel the production of documents by the same means and, so far as may be, in the same manner as is provided in the case of a civil court under the Code of Civil Procedure, 1908 (V of 1908).



(2) The provisions of the Indian Evidence Act, 1872, shall, subject to the provisions of these rules and the Standing Orders made by the President, be deemed to apply to every such inquiry.



60. (1) If in the opinion of the Credentials Committee or of the election tribunal, as the case may be, the election of the returned candidate has been vitiated by a corrupt practice of the kind specified in the Indian Legislative Assembly Electoral Rules as in force on the 1st day of August, 1947 or has been materially affected by the improper acceptance or refusal of any nomination or by the improper reception or refusal of a vote or the reception of any vote which is void, or if the election has not been a free election by reason of the large number of cases in which undue influence or bribery of the kind described in the aforesaid Legislative Assembly Electoral Rules has been exercised or committed, the Committee or the tribunal, as the case may be may recommend in the report that the election be declared void.



(2) The report of the Credentials Committee, or the election tribunal, as the case may be, shall include a recommendation as to the total amount of costs which are payable and the persons by and to whom such costs should be paid out of the sum deposited as security under rule 55 and whether the said sum should be returned.



61. On receipt of the report of the Credentials Committee or the election tribunal, as the case may be, the President shall issue orders in accordance therewith and the orders so issued shall be final and shall not be questioned in any court.



\begin{center}
CHAPTER XII\end{center}
\begin{center}
\emph{Miscellaneous}\end{center}


62. All elections in the Assembly to be held on the principle of proportional representation by means of the single transferable vote shall be conducted \emph{mutatis mutandis} in accordance with the regulations in force in this behalf in the Indian Legislative Assembly.



63. Subject to the requirement of a quorum prescribed by or under these rules, the Assembly and any committee set up by the Assembly shall have power to act notwithstanding any vacancy in the membership thereof.



64. No new rule shall be made nor shall any of these rules be amended or deleted except after a reference of the proposal so to make, amend, or delete the rule to the Steering Committee, which shall report to the Assembly within two weeks of the receipt of the reference.



65. Save as otherwise provided in these rules, the provisions thereof shall apply \emph{mutatis mutandis} to the Committees of the Assembly.



66. Where, in the opinion of the President, any difficulty arises in the carrying out of these rules, or in respect of any matter for which no provision is made in these rules, the President may, notwithstanding anything contained therein, make such provision as he thinks fit for the purpose of removing the difficulty.



67. If any question arises as to the interpretation of these rules otherwise than in connection with an election held thereunder, the question shall be referred for the decision of the President and his decision shall be final.



\begin{center}
THE SCHEDULE \end{center}
\begin{center}
(See Rule 51) \end{center}


1. The seats allotted to Indian States in the Statement shall be allocated among the various States and groups of States as in Annexure A, generally on the basis of one seat for one million of the population, fractions of three-fourths or more being counted as one and lesser fraction being ignored in the case of individual States, and fractions of more than half being counted as one and lesser fractions being ignored in the case of groups of States.



2. The President may, on the application of any State or States concerned, by order amend Annexure A to this Schedule so as to—



(a) alter the representation allotted to the States, individual or grouped;



(b) alter the grouping of the States by the division of a group into more than one group or the transfer of any State, or States from one group to another or otherwise; 



Provided that— 



(i) no such alteration shall affect the total representation of all States or of the group or groups of States concerned; and 



(ii) in making any such alteration the population basis shall not be departed from and the geographical proximity, economic considerations, and ethnic, cultural and linguistic affinity shall be duly kept in view.



2-A. When the representation allotted to the States, individual or grouped, or the grouping of the States is altered by an order made under paragraph 2, the President may, on application made in that behalf by the States affected by such order, declare the seats of the members of the Assembly representing the States so affected to be vacant.



3. Not less than 50 per cent of the total representatives of the States in the Assembly shall be elected by the elected members of the States’ legislatures, or where, such legislatures do not exist, by the members of electoral colleges constituted in accordance with the provisions made in this behalf by the Rulers of the States concerned. The States shall endeavour to increase the quota of elected representatives as much above 50 per cent of the total number as possible. Accordingly at least one half of the number of seats allotted to any State or group of States shall be filled by election in accordance with the provision made in that behalf by the Ruler of the State or States concerned. 



4. The Conveners, in respect of the various groups of States specified in column 1 of the Annexure A, shall be the rulers specified in the corresponding entries in column 4 of that Annexure. The Secretary may in consultation with the States in the group make any such changes in the said column 4 as he may deem necessary or desirable. 



5. On the completion of the election or nomination, as the case may be, the Ruler of the State concerned shall make a notification as far as may be in the following form stating the name or names of the person or persons elected or nominated as representative or representatives in the Constituent Assembly and cause it to be communicated to the President of the Constituent Assembly. Where the selection has been made by a group of States, this notification shall be made by the convener for that group.



\begin{center}
FORM \end{center}
BE IT HEREBY KNOWN THAT [here enter the name of the representative(s)] \_\_\_\_\_ has/have been duly chosen as (a) representative(s) of [here, enter the name(s) of the State(s)] \_\_\_\_\_ in the Constituent Assembly of India. In testimony whereof this notification is issued under my signature and the Seal of my State. 

State(s) \_\_\_\_\_\_\_\_\_\_

Date \_\_\_\_\_\_\_\_\_\_

\begin{flushright}
Ruler of \_\_\_\_\_\_\_\_\_\_\end{flushright}
\begin{center}
ANNEXURE A\end{center}
\begin{center}
\emph{Single State}\end{center}


\begin{tabular}{|c|c|c|c|} \hlineDivision as shown in the Table of Seats appended to Part II of the First Schedule to the Govt. of India Act, 1935.

Name of State

Number of seats in the Constituent Assembly

Convener

\\ \hline1

2

3

4

\\ \hlineI

Hyderabad

16

..

\\ \hlineII

Mysore

7

..

\\ \hlineIII

Kashmir

4

..

\\ \hlineIV

Gwalior

4

..

\\ \hlineV

Baroda

3

..

\\ \hlineIX

Travancore

6

..

\\ \hlineIX

Cochin

1

..

\\ \hlineX

Udaipur

2

..

\\ \hlineX

Jaipur

2

..

\\ \hlineX

Jodhpur

2

..

\\ \hlineX

Bikaner

1

..

\\ \hlineX

Alwar

1

..

\\ \hlineX

Kotah

1

..

\\ \hlineXI

Indore

1

..

\\ \hlineXI

Bhopal

1

..

\\ \hlineXI

Rewa

1

..

\\ \hlineXII

Kolhapur

1



\\ \hlineXIV

Patiala

2



\\ \hlineXIV

Bahawalpur

1



\\ \hlineXVI

Mayurbhanj

1



\\ \hline

20

60



\\ \hline\end{tabular}

\begin{center}
Frontier Groups	\end{center}


\begin{tabular}{|c|c|c|c|} \hlineVII

Sikkim



1



Ruler of:—

Cooch Behar State.

\\ \hlineXV

Cooch Behar

\\ \hlineXV

Tripura





1





Tripura State.

\\ \hlineXV

Manipur

\\ \hlineXVII

Khasi States

\\ \hline\end{tabular}\begin{center}
\end{center}
\begin{center}
Interior Groups\end{center}


\begin{tabular}{|c|c|c|c|} \hlineVIII



Rampur



1



Rampur State

\\ \hlineBenares

\\ \hlineX









(13 States)











Bharatpur







3



3

























Bundi State















\\ \hlineTonk

\\ \hlineDholpur

\\ \hlineKarauli

\\ \hlineBundi

\\ \hlineSirohi

\\ \hlineDungarpur

\\ \hlineBanswara

\\ \hlinePartabgarh

\\ \hlineJhalawar

\\ \hlineJaisalmer

\\ \hlineKishengarh

\\ \hlineXI

Shahpura

\\ \hlineXVII

























(26 States)























Datia



























3



















































Panna State.

























\\ \hlineOrcha

\\ \hlineDhar

\\ \hlineDewas (Senior)

\\ \hlineDewas (Junior

\\ \hlineJaora

\\ \hlineRatlam

\\ \hlinePanna

\\ \hlineSamthar

\\ \hlineAjaigarh

\\ \hlineBijawar

\\ \hlineCharkhari

\\ \hlineChhatarpur

\\ \hlineBaoni

\\ \hlineNagod

\\ \hlineMaihar

\\ \hlineBaraundha

\\ \hlineBarwani

\\ \hlineAli Rajpur

\\ \hlineJhabua

\\ \hlineSailana

\\ \hlineSitamau

\\ \hlineRaigarh

\\ \hlineNarsingarh 

\\ \hlineKhilchipur

\\ \hlineXVII

Kurwai

\\ \hlineXII











(17 States)





















Cutch













4

































Nawanagar State.





















\\ \hlineIdar

\\ \hlineNawanagar

\\ \hlineBhavnagar

\\ \hlineJunagadh

\\ \hlineDhrangadhra

\\ \hlineGondal

\\ \hlinePorbandar

\\ \hlineMorvi

\\ \hlineRadhanpur

\\ \hlineWankaner

\\ \hlinePalitana

\\ \hlineDhrol

\\ \hlineLimbdi

\\ \hlineWadhwan

\\ \hlineRajkot

\\ \hlineJafrabad

\\ \hlineXII-A













(14 States)











Rajpipla















2



























Rajpipla State.













\\ \hlinePalanpur

\\ \hlineCambay

\\ \hlineDharampur

\\ \hlineBalasinor

\\ \hlineBaria

\\ \hlineChhota Udepur

\\ \hlineSant

\\ \hlineLunawada

\\ \hlineBansda

\\ \hlineSachin

\\ \hlineJawhar

\\ \hlineDanta

\\ \hlineXIII

Janjira

\\ \hlineXIII











(17 States)















Sangh













2

































Miraj (Junior) State.





















\\ \hlineSavantvadi

\\ \hlineMudhol

\\ \hlineBhor

\\ \hlineJamkhandi

\\ \hlineMiraj (Senior)

\\ \hlineMiraj (Junior)

\\ \hlineKurundwad (Senior)

\\ \hlineKurundwad (Junior)

\\ \hlineAkalkot

\\ \hlinePhaltan

\\ \hlineJath

\\ \hlineAundh

\\ \hlineRamdurg

\\ \hlineIX





Pudukkottai

\\ \hlineBanganapallee

\\ \hlineSandur

\\ \hlineXIV











(14 States)













Kapurthala













3



























Bilaspur State.















\\ \hlineJind

\\ \hlineNabha

\\ \hlineMandi

\\ \hlineBilaspur

\\ \hlineSuket

\\ \hlineTehri-Garhwal

\\ \hlineSirmur

\\ \hlineChamba

\\ \hlineFaridkot

\\ \hlineMalerkotla

\\ \hline*Loharu

\\ \hlineKalsia

\\ \hlineXVII

Bashahr

\\ \hlineXV











(25 States)













Sonepur











4

















































Bundi State.







































\\ \hlinePatna

\\ \hlineKalahandi

\\ \hlineKeonjhar

\\ \hlineDhenkanal

\\ \hlineNayagarh

\\ \hlineTalcher

\\ \hlineNilgiri

\\ \hlineGangpur

\\ \hlineBamra

\\ \hlineSeraikela

\\ \hlineBaud

\\ \hlineBonai

\\ \hlineXVII























Athgarh

\\ \hlinePal Lahara

\\ \hlineAthmalik

\\ \hlineHindol

\\ \hlineNarsingpur

\\ \hlineBaramba

\\ \hlineTigiria

\\ \hlineKhandpara

\\ \hlineRanpur

\\ \hlineDaspalla

\\ \hlineRairakhol

\\ \hlineKharsawan

\\ \hlineXVI-A









(14 States)







Bastar











3



























Baud State.

















\\ \hlineSurguja

\\ \hlineRaigarh

\\ \hlineNandgaon

\\ \hlineKhairagarh

\\ \hlineJaipur

\\ \hlineKanker

\\ \hlineKorea

\\ \hlineSarangarh

\\ \hlineXVII









Changbhakar

\\ \hlineChhuikadan

\\ \hlineKawardha

\\ \hlineSakti

\\ \hlineUdaipur

\\ \hlineXVII

All other States

4

Baghat State.

\\ \hline\end{tabular}

*By special arrangement Loharu is represented by the representative of Bikaner State.

